%! Bias in Samples and Surveys — Unit 1, Lesson 1.4 — Student Packet Cover Sheet
% PILOT SHAPE (2026-09-06): modeled on AP Statistics lesson 1.4 minus its AP Practice page.
% THREE packet rows — Warm-Up, Guided Notes & Practice, Homework. The homework is the individual
% practice, started in class, and it is SCORED, so its cell is a \blank{}, never NA.
\documentclass[12pt]{article}
\usepackage{saar-article}
\usepackage{saar-boxes}
\usepackage{ltablex}
\keepXColumns
% Measured banner: the band grows to fit the title block, so a title long enough to wrap grows
% the band rather than running out of it (the stats course's \coverbanner; saar-article does not
% carry one, so it is defined here — every cover copies this block verbatim).
\newsavebox{\bannerbox}
\newlength{\bannerht}
\newlength{\coverbannerpad}
\setlength{\coverbannerpad}{0.16in}       % breathing room above and below the text
\newcommand{\coverbanner}[2]{%
  \sbox{\bannerbox}{%
    \begin{minipage}{\dimexpr\paperwidth-1.4in\relax}%
      \centering
      {\color{white}\textbf{\LARGE \CourseName}}\\[4pt]
      {\color{frostmid}\large\textbf{#1}}\\[6pt]
      {\color{white}\textbf{\large #2}}%
    \end{minipage}}%
  \setlength{\bannerht}{\dimexpr\ht\bannerbox+\dp\bannerbox+2\coverbannerpad\relax}%
  \noindent\begin{tikzpicture}[remember picture, overlay]
    \fill[cerulean] (current page.north west)
      rectangle ([yshift=-\bannerht]current page.north east);
    \node[anchor=north, inner sep=0pt]
      at ([yshift=-\coverbannerpad]current page.north) {\usebox{\bannerbox}};
  \end{tikzpicture}\par
  % The text body starts 0.6in below the page edge (geometry top=0.6in), so reserve only what
  % the band overruns that, plus a gap under the band.
  \vspace*{\dimexpr\bannerht-0.6in+0.16in\relax}%
}

\begin{document}

% Full-bleed cerulean banner, sized to the title block.
\coverbanner{Unit 1}{Lesson 1.4 \quad Bias in Samples and Surveys}

\vspace{0.06in}
% NAMESTRIP: this is the ONLY component that carries the name row.
\namedateperiod
\vspace{0.18in}

% GRADUAL RELEASE: the vocabulary is taught in the guided notes, so the targets name it
% plainly. The standard codes (PS.DC.2c; AFDA.DA.2e-g) stay in the teacher-facing plan.
\begin{learningtargetbox}
\small
\begin{enumerate}[leftmargin=1.5em, itemsep=5pt]
  \item \ldots{} explain what makes a study \textbf{biased} and tell bias apart from
        \textbf{sampling variability} --- a direction, not bad luck.
  \item \ldots{} identify \textbf{sampling bias} in a plan, naming \textbf{undercoverage} and
        \textbf{nonresponse} --- by asking \emph{could that person have been chosen?}
  \item \ldots{} name the forms of \textbf{response bias}, including \textbf{question order
        bias}, and say which \textbf{direction} a biased plan pushes the estimate.
  \item \ldots{} explain why a \textbf{larger sample does not fix bias}, and describe and
        \emph{justify} a change that removes it.
\end{enumerate}
\end{learningtargetbox}

\vspace{0.14in}

\begin{tocbox}
\small
\renewcommand{\arraystretch}{1.5}
\begin{tabularx}{\linewidth}{@{} c l X r @{}}
\rowcolor{frost}
\textbf{\#} & \textbf{Component} & \textbf{Description} & \textbf{Score} \\
\midrule
1 & Warm-Up & Harbor Point's $1{,}500$ members: a percent, a sample's percent, and a
             prediction set beside the truth & \blank{1.2cm} \\
2 & Guided Notes \& Practice & Bias as a direction; who never had a chance; response bias;
             the ten-week fix that fixes nothing --- then Cedar Ridge Apartments worked
             together & \blank{1.2cm} \\
% Row 3 is the lesson's GRADED practice, started in class. Packet pages are the default; a
% Desmos activity may replace them on a given day, decided at Close & Assign. The row and its
% score cell never change.
3 & Homework & Millbrook Public Library: two studies, two answers --- \textbf{scored, started
             in class, due the first class after two study halls} & \blank{1.2cm} \\
\midrule
  & \multicolumn{2}{r}{\textbf{Total}} & \blank{1.2cm} \\
\end{tabularx}
\end{tocbox}

\vspace{0.14in}

% Keep in Mind carries the lesson's CONTENT — the rule, the test, the trap — never its process.
\begin{remindbox}
\small
\textbf{Bias is a direction, not bad luck.} Honest random samples scatter \emph{around} the
truth; a biased procedure misses the \emph{same way every time}. Before anyone speaks:
\textbf{undercoverage} (left off the frame) and \textbf{nonresponse} (chosen but silent) ---
ask \emph{could that person have been chosen?} After the right person is chosen:
\textbf{response bias} --- demand, social desirability, dissent, acquiescence, extreme
responses, neutral responding, question order. \textbf{Watch out:} a larger sample makes an
estimate more \emph{precise}, not more \emph{correct} --- size shrinks variability and does
nothing to direction. Name the bias, say which way it pushes, and attach the fix to what it
removes.
\end{remindbox}

\end{document}
